\chapter{2012年天津卷（理）}

\section{单选题}

\begin{problem}
\(i\) 是虚数单位，复数 \(\frac{7-\i}{3+\i}=\)（\quad）

\choices
    {\(2 + \i\)}
    {\(2 - \i\)}
    {\(-2 + \i\)}
    {\(-2 - \i\)}
\end{problem}

\begin{answer}
B
\end{answer}

\begin{solution}
分子、分母同乘以分母的共轭复数 \(3-\i\)，得
\[
\frac{7-\i}{3+\i}=\frac{(7-\i)(3-\i)}{(3+\i)(3-\i)}=\frac{20-10\i}{10}=2-\i
\]
因此选 B.
\end{solution}

\begin{problem}
设 \(\varphi \in \R\)，则“\(\varphi = 0\)”是“\(f(x) = \cos (x + \varphi)\)（\(x \in \R\)）为偶函数”的（\quad）

\choices
    {充分而不必要条件}
    {必要而不充分条件}
    {充分必要条件}
    {既不充分也不必要条件}
\end{problem}

\begin{answer}
A
\end{answer}

\begin{solution}
由和角公式，\(f(x)=\cos x\cos\varphi-\sin x\sin\varphi\)，从而 \(f(-x)=\cos x\cos\varphi+\sin x\sin\varphi\). 因此 \(f(x)-f(-x)=-2\sin x\sin\varphi\). 要使 \(f(x)\) 对任意 \(x\in\R\) 为偶函数，必须且只须 \(\sin\varphi=0\)，即 \(\varphi=k\pi\)（\(k\in\Z\)）. 当 \(\varphi=0\) 时确有 \(f(x)=\cos x\) 为偶函数；但取 \(\varphi=\pi\) 时 \(f(x)=-\cos x\) 也为偶函数，故 \(\varphi=0\) 不是必要条件. 因此选 A.
\end{solution}

\begin{problem}
阅读如图所示的程序框图，运行相应的程序，当输入 \(x\) 的值为 \(-25\) 时，输出 \(x\) 的值为（\quad）

\bitmapfigure[max width=0.42\linewidth]{img/2012/tianjin_science/q3_fig1.png}

\choices
    {\(-1\)}
    {\(1\)}
    {\(3\)}
    {\(9\)}
\end{problem}

\begin{answer}
C
\end{answer}

\begin{solution}
程序先判断 \(|x|>1\)，若成立则把 \(x\) 替换为 \(\sqrt{|x|}-1\). 从 \(x=-25\) 开始，第一次循环后 \(x=\sqrt{25}-1=4\)，第二次循环后 \(x=\sqrt4-1=1\). 此时 \(|x|>1\) 不成立，执行 \(x=2x+1\)，得到输出值 \(3\). 因此选 C.
\end{solution}

\begin{problem}
函数 \( f(x) = 2^x + x^3 - 2 \) 在区间 \((0,1)\) 内的零点个数是（\quad）

\choices
    {\(0\)}
    {\(1\)}
    {\(2\)}
    {\(3\)}
\end{problem}

\begin{answer}
B
\end{answer}

\begin{solution}
函数在 \([0,1]\) 上连续，且 \(f(0)=1+0-2=-1<0\)，\(f(1)=2+1-2=1>0\)，由零点存在定理知它在 \((0,1)\) 内至少有一个零点. 又因为在 \((0,1)\) 内
\[
f'(x)=2^x\ln2+3x^2>0
\]
所以 \(f(x)\) 在该区间上严格递增，至多有一个零点. 综上，零点个数为 \(1\)，因此选 B.
\end{solution}

\begin{problem}
在 \(\left(2x^{2} - \frac{1}{x}\right)^{5}\) 的二项展开式中，\(x\) 的系数为（\quad）

\choices
    {\(10\)}
    {\(-10\)}
    {\(40\)}
    {\(-40\)}
\end{problem}

\begin{answer}
D
\end{answer}

\begin{solution}
二项展开式的第 \(k+1\) 项为
\[
\mathrm C_5^k(2x^2)^{5-k}\left(-\frac1x\right)^k=\mathrm C_5^k2^{5-k}(-1)^k x^{10-3k}
\]
要得到 \(x\) 项，须有 \(10-3k=1\)，解得 \(k=3\). 因此所求系数为
\[
\mathrm C_5^3 2^2(-1)^3=-40
\]
所以选 D.
\end{solution}

\begin{problem}
在 \(\triangle ABC\) 中，内角 \(A\)，\(B\)，\(C\) 所对的边分别是 \(a\)，\(b\)，\(c\)，已知 \(8b = 5c\)，\(C = 2B\)，则 \(\cos C =\)（\quad）

\choices
    {\(\frac{7}{25}\)}
    {\( -\frac{7}{25} \)}
    {\(\pm\frac{7}{25}\)}
    {\(\frac{24}{25}\)}
\end{problem}

\begin{answer}
A
\end{answer}

\begin{solution}
由三角形内角的范围及 \(C=2B\)，有 \(0<B<\frac\pi2\)，所以 \(\sin B>0\) 且 \(\cos B>0\). 由正弦定理，\(\frac{b}{c}=\frac{\sin B}{\sin C}=\frac{5}{8}\). 又 \(C=2B\)，故
\[
\frac{\sin B}{\sin2B}=\frac{\sin B}{2\sin B\cos B}=\frac{1}{2\cos B}=\frac58
\]
从而 \(\cos B=\frac45\). 因此
\[
\cos C=\cos2B=2\cos^2B-1=2\left(\frac45\right)^2-1=\frac7{25}
\]
所以选 A.
\end{solution}

\begin{problem}
已知 \(\triangle ABC\) 为等边三角形，\(AB = 2\)，设点 \(P\)，\(Q\) 满足 \(\overrightarrow{AP} = \lambda \overrightarrow{AB}\)，\(\overrightarrow{AQ} = (1 - \lambda) \overrightarrow{AC}\)，\(\lambda \in \R\)，若 \(\overrightarrow{BQ} \cdot \overrightarrow{CP} = -\frac{3}{2}\)，则 \(\lambda =\)（\quad）

\choices
    {\( \frac{1}{2} \)}
    {\( \frac{1 \pm \sqrt{2}}{2} \)}
    {\(\frac{1 \pm \sqrt{10}}{2}\)}
    {\(\frac{-3 \pm 2\sqrt{2}}{2}\)}
\end{problem}

\begin{answer}
A
\end{answer}

\begin{solution}
以 \(A\) 为原点，令 \(\overrightarrow{AB}=(2,0)\)，\(\overrightarrow{AC}=(1,\sqrt3)\). 于是 \(\overrightarrow{AP}=(2\lambda,0)\)，\(\overrightarrow{AQ}=((1-\lambda),(1-\lambda)\sqrt3)\).
从而
\(\overrightarrow{BQ}=(-1-\lambda,(1-\lambda)\sqrt3)\)，\(\overrightarrow{CP}=(2\lambda-1,-\sqrt3)\).
两向量的数量积为
\[
\overrightarrow{BQ}\cdot\overrightarrow{CP}=(-1-\lambda)(2\lambda-1)-3(1-\lambda)=-2\lambda^2+2\lambda-2
\]
由题设得 \(-2\lambda^2+2\lambda-2=-\frac32\)，即 \(\lambda^2-\lambda+\frac14=0\)，所以 \(\lambda=\frac12\). 因此选 A.
\end{solution}

\begin{problem}
设 \(m\)，\(n \in \R\)，若直线 \((m + 1)x + (n + 1)y - 2 = 0\) 与圆 \((x - 1)^2 + (y - 1)^2 = 1\) 相切，则 \(m + n\) 的取值范围是（\quad）

\choices
    {\([1 - \sqrt{3}, 1 + \sqrt{3}]\)}
    {\((-\infty, 1 - \sqrt{3}] \cup [1 + \sqrt{3}, +\infty)\)}
    {\([2 - 2\sqrt{2}, 2 + 2\sqrt{2}]\)}
    {\((-\infty, 2 - 2\sqrt{2}] \cup [2 + 2\sqrt{2}, +\infty)\)}
\end{problem}

\begin{answer}
D
\end{answer}

\begin{solution}
圆心为 \((1,1)\)，半径为 \(1\). 由点到直线的距离等于半径，得
\[
\frac{|(m+1)+(n+1)-2|}{\sqrt{(m+1)^2+(n+1)^2}}=1
\]
即 \((m+n)^2=(m+1)^2+(n+1)^2\)，化简为 \((m-1)(n-1)=2\). 令 \(u=m-1\)，\(v=n-1\)，则 \(uv=2\). 当 \(u\)，\(v>0\) 时，\(u+v\geqslant2\sqrt2\)；当 \(u\)，\(v<0\) 时，\(u+v\leqslant-2\sqrt2\). 由于 \(m+n=u+v+2\)，所以
\[
m+n\in(-\infty,2-2\sqrt2]\cup[2+2\sqrt2,+\infty)
\]
因此选 D.
\end{solution}

\section{填空题}

\begin{problem}
某地区有小学 150 所，中学 75 所，大学 25 所. 现采用分层抽样的方法从这些学校中抽取 30 所学校对学生进行视力调查，应从小学中抽取\fillinblank{} 所学校，中学中抽取\fillinblank{} 所学校.
\end{problem}

\begin{answer}
\(18\)，\(9\)
\end{answer}

\begin{solution}
学校总数为 \(150+75+25=250\) 所，抽样比例为 \(\frac{30}{250}=\frac3{25}\). 因此小学抽取 \(150\times\frac3{25}=18\) 所，中学抽取 \(75\times\frac3{25}=9\) 所.
\end{solution}

\begin{problem}
一个几何体的三视图如图所示（单位：\(\mathrm{m}\)），则该几何体的体积为\fillinblank{}\(\mathrm{m}^3\).

\bitmapfigure[max width=0.74\linewidth]{img/2012/tianjin_science/q10_fig1.png}
\end{problem}

\begin{answer}
\(18 + 9\pi\)
\end{answer}

\begin{solution}
由三视图可知，几何体由一个长为 \(6\)，宽为 \(3\)，高为 \(1\) 的长方体和两个半径为 \(\frac32\) 的球组成. 长方体体积为 \(6\times3\times1=18\)，两个球的体积为
\[
2\times\frac43\pi\left(\frac32\right)^3=9\pi
\]
所以几何体的体积为 \(18+9\pi\mathrm{~m}^3\).
\end{solution}

\begin{problem}
已知集合 \(A = \{x \in \R \mid |x + 2| < 3\}\)，集合 \(B = \{x \in \R \mid (x - m)(x - 2) < 0\}\)，且 \(A \cap B = (-1, n)\)，则 \(m =\fillinblank{}\)，\(n =\fillinblank{}\).
\end{problem}

\begin{answer}
\(-1\)，\(1\)
\end{answer}

\begin{solution}
由 \(|x+2|<3\)，得 \(-5<x<1\)，即 \(A=(-5,1)\). 若 \(m\neq2\)，则 \(B\) 是 \(m\) 与 \(2\) 之间的开区间；若 \(m=2\)，则 \(B\) 为空集，与题设不符. 交集的左端点为 \(-1\)，而 \(A\) 的左端点为 \(-5\)，所以必须有 \(\min\{m,2\}=-1\)，从而 \(m=-1\). 此时 \(B=(-1,2)\)，故 \(A\cap B=(-1,1)\)，所以 \(n=1\).
\end{solution}

\begin{problem}
已知抛物线的参数方程为 \(\begin{cases}
x = 2pt^2,\\
y = 2pt,
\end{cases}\)（\(t\) 为参数），其中 \(p > 0\)，焦点为 \(F\)，准线为 \(l\)，过抛物线上一点 \(M\) 作 \(l\) 的垂线，垂足为 \(E\)，若 \(|EF| = |MF|\)，点 \(M\) 的横坐标是 3，则 \(p = \)\fillinblank{}.
\end{problem}

\begin{answer}
2
\end{answer}

\begin{solution}
消去参数得 \(y^2=2px\)，所以抛物线的焦点为 \(F\left(\frac p2,0\right)\)，准线为 \(l:x=-\frac p2\). 设点 \(M\) 对应的参数为 \(t\)，则 \(M=(2pt^2,2pt)\)，\(E=\left(-\frac p2,2pt\right)\).
因此
\[
|EF|=p\sqrt{1+4t^2}
\]
并且
\[
|MF|^2=\left(2pt^2-\frac p2\right)^2+(2pt)^2=p^2\left(2t^2+\frac12\right)^2
\]
由 \(p>0\) 得 \(|MF|=p\left(2t^2+\frac12\right)\). 由 \(|EF|=|MF|\)，两边除以 \(p\) 后平方，得
\[
1+4t^2=\left(2t^2+\frac12\right)^2
\]
即 \(4t^4-2t^2-\frac34=0\)，所以 \(t^2=\frac34\)（另一根为负数，舍去）. 又因点 \(M\) 的横坐标为 \(2pt^2=3\)，故 \(\frac32p=3\)，从而 \(p=2\).
\end{solution}

\begin{problem}
如图，已知 \(AB\) 和 \(AC\) 是圆的两条弦，过点 \(B\) 作圆的切线与 \(AC\) 的延长线相交于点 \(D\). 过点 \(C\) 作 \(BD\) 的平行线与圆相交于点 \(E\)，与 \(AB\) 相交于点 \(F\)，\(AF = 3\)，\(FB = 1\)，\(EF = \frac{3}{2}\)，则线段 \(CD\) 的长为\fillinblank{}.

\bitmapfigure[max width=0.58\linewidth]{img/2012/tianjin_science/q13_fig1.png}
\end{problem}

\begin{answer}
\(\frac{4}{3}\)
\end{answer}

\begin{solution}
由点 \(F\) 向圆作两条割线，利用相交弦定理得
\[
FA\cdot FB=FC\cdot FE
\]
所以 \(FC=\frac{3\times1}{3/2}=2\). 又因为 \(CF\parallel BD\)，所以 \(\triangle ACF\sim\triangle ADB\)，从而
\[
\frac{AC}{AD}=\frac{AF}{AB}=\frac34
\]
设 \(CD=x\)，则 \(AC=3x\)，\(AD=4x\). 由点 \(D\) 的切线割线定理，得
\[
DB^2=DC\cdot DA=4x^2
\]
即 \(DB=2x\). 相似三角形还给出 \(\frac{CF}{DB}=\frac{AF}{AB}=\frac34\)，故 \(DB=\frac43CF=\frac83\). 因此 \(2x=\frac83\)，解得 \(CD=x=\frac43\).
\end{solution}

\begin{problem}
已知函数 \( y = \frac{|x^2 - 1|}{x - 1} \) 的图象与函数 \( y = kx - 2 \) 的图象恰有两个交点，则实数 \( k \) 的取值范围是\fillinblank{}.
\end{problem}

\begin{answer}
\((0,1)\cup(1,4)\)
\end{answer}

\begin{solution}
函数定义域为 \(x\neq1\). 因 \(|x^2-1|=|x-1||x+1|\)，所以
\[
\frac{|x^2-1|}{x-1}=\begin{cases}
x+1,&x\leqslant-1\text{ 或 }x>1,\\
-x-1,&-1<x<1.
\end{cases}
\]
与直线 \(y=kx-2\) 联立第一段，若 \(k\neq1\)，交点横坐标为 \(x_1=\frac3{k-1}\). 要求 \(x_1\leqslant-1\) 或 \(x_1>1\)，得到 \(k\in[-2,1)\cup(1,4)\)；\(k=1\) 时两直线平行，无交点.

联立第二段，若 \(k\neq-1\)，交点横坐标为 \(x_2=\frac1{k+1}\). 要求 \(-1<x_2<1\)，得到 \(k\in(-\infty,-2)\cup(0,+\infty)\)；\(k=-1\) 时两直线平行. 因此两段各有一个交点的参数必须属于
\[
\bigl([-2,1)\cup(1,4)\bigr)\cap\bigl((-\infty,-2)\cup(0,+\infty)\bigr)=(0,1)\cup(1,4)
\]
这正是图象共有两个交点的情形；其余取值下交点数不是两个. 因此
\[
k\in(0,1)\cup(1,4)
\]
\end{solution}

\section{解答题}

\begin{problem}
已知函数 \(f(x) = \sin \left(2x + \frac{\pi}{3}\right) + \sin \left(2x - \frac{\pi}{3}\right) + 2\cos^2 x - 1\)，\(x \in \R\).

\begin{enumerate}
\item 求函数 \( f(x) \) 的最小正周期；
\item 求函数 \( f(x) \) 在区间 \( \left[-\frac{\pi}{4}, \frac{\pi}{4}\right] \) 上的最大值和最小值.
\end{enumerate}

\end{problem}

\begin{solution}
由和差化积公式及二倍角公式，
\[
f(x)=2\sin2x\cos\frac\pi3+\cos2x=\sin2x+\cos2x=\sqrt2\sin\left(2x+\frac\pi4\right)
\]
\begin{enumerate}
\item 函数 \(\sin(2x+\frac\pi4)\) 的最小正周期为 \(\frac{2\pi}{2}=\pi\)，故 \(f(x)\) 的最小正周期为 \(\pi\).
\item 令 \(t=2x\)，则 \(t\in[-\frac\pi2,\frac\pi2]\)，且 \(f(x)=\sin t+\cos t\). 其导数为 \(\cos t-\sin t\)，在该区间内于 \(t=\frac\pi4\) 处由正变负，所以此时取得最大值 \(\sqrt2\). 端点处 \(f(-\frac\pi4)=-1\)，\(f(\frac\pi4)=1\)，故最小值为 \(-1\).
\end{enumerate}
\end{solution}

\begin{problem}
有4个人去参加某娱乐活动，该活动有甲，乙两个游戏可供参加者选择. 为增加趣味性，约定：每个人通过掷一枚质地均匀的骰子决定自己去参加哪个游戏，掷出点数为1或2的人去参加甲游戏，掷出点数大于2的人去参加乙游戏.

\begin{enumerate}
\item 求这 4 个人中恰有 2 人去参加甲游戏的概率；
\item 求这 4 个人中去参加甲游戏的人数大于去参加乙游戏的人数的概率；
\item 用 \(X\)，\(Y\) 分别表示这 4 个人中去参加甲，乙游戏的人数，记 \( \xi = |X - Y| \)，求随机变量 \( \xi \) 的分布列与数学期望 \( E(\xi) \).
\end{enumerate}

\end{problem}

\begin{solution}
每个人参加甲游戏的概率为 \(\frac13\)，参加乙游戏的概率为 \(\frac23\). 设恰有 \(i\) 人参加甲游戏的事件为 \(A_i\)，则
\[
P(A_i)=\mathrm C_4^i\left(\frac13\right)^i\left(\frac23\right)^{4-i}
\]
\begin{enumerate}
\item
\[
P(A_2)=\mathrm C_4^2\left(\frac13\right)^2\left(\frac23\right)^2=\frac8{27}
\]
\item 甲游戏人数多于乙游戏人数等价于 \(X>2\)，所以
\[
P(X>2)=P(A_3)+P(A_4)=\mathrm C_4^3\left(\frac13\right)^3\frac23+\left(\frac13\right)^4=\frac19
\]
\item 因为 \(X+Y=4\)，所以 \(\xi=|2X-4|\) 的可能取值为 \(0\)，\(2\)，\(4\). 于是
\[
P(\xi=0)=P(A_2)=\frac8{27}
\]
\[
P(\xi=2)=P(A_1)+P(A_3)=\frac{32}{81}+\frac8{81}=\frac{40}{81}
\]
\[
P(\xi=4)=P(A_0)+P(A_4)=\frac{16}{81}+\frac1{81}=\frac{17}{81}
\]
故分布列为
\begin{center}
\begin{tabular}{c|ccc}
\(\xi\)&\(0\)&\(2\)&\(4\)\\
\hline
\(P\)&\(\frac8{27}\)&\(\frac{40}{81}\)&\(\frac{17}{81}\)
\end{tabular}
\end{center}
数学期望为
\[
E(\xi)=0\times\frac8{27}+2\times\frac{40}{81}+4\times\frac{17}{81}=\frac{148}{81}
\]
\end{enumerate}
\end{solution}

\begin{problem}
如图，在四棱锥 \(P - ABCD\) 中，\(PA \perp\) 平面 \(ABCD\)，\(AC \perp AD\)，\(AB \perp BC\)，\(\angle BAC = 45^{\circ}\)，\(PA = AD = 2\)，\(AC = 1\).

\begin{enumerate}
\item 证明 \(PC \perp AD\)；
\item 求二面角 \(A - PC - D\) 的正弦值；
\item 设 \(E\) 为棱 \(PA\) 上的点，满足异面直线 \(BE\) 与 \(CD\) 所成的角为 \(30^{\circ}\)，求 \(AE\) 的长.
\end{enumerate}
\bitmapfigure[max width=0.42\linewidth]{img/2012/tianjin_science/q17_fig1.png}

\end{problem}

\begin{solution}
建立空间直角坐标系，使 \(A\) 为原点，\(AD\) 为 \(x\) 轴，\(AC\) 为 \(y\) 轴，\(AP\) 为 \(z\) 轴. 于是 \(D=(2,0,0)\)，\(C=(0,1,0)\)，\(P=(0,0,2)\). 设 \(B=(u,v,0)\). 由 \(\angle BAC=45^{\circ}\) 且 \(v>0\)，有
\[
\frac{v}{\sqrt{u^2+v^2}}=\frac{\sqrt2}{2}
\]
从而 \(u^2=v^2\). 又由 \(AB\perp BC\)，得
\[
(u,v,0)\cdot(-u,1-v,0)=0
\]
即 \(u^2+v^2=v\). 因 \(v>0\)，解得 \(v=\frac12\)，\(u=\pm\frac12\). 由四边形 \(ABCD\) 的顶点顺序，\(B\) 与 \(D\) 位于直线 \(AC\) 的两侧，故 \(u=-\frac12\). 因此 \(A=(0,0,0)\)，\(D=(2,0,0)\)，\(C=(0,1,0)\)，\(B=\left(-\frac12,\frac12,0\right)\)，\(P=(0,0,2)\).
\begin{enumerate}
\item \(\overrightarrow{PC}=(0,1,-2)\)，\(\overrightarrow{AD}=(2,0,0)\)，所以 \(\overrightarrow{PC}\cdot\overrightarrow{AD}=0\)，从而 \(PC\perp AD\).
\item 平面 \(PAC\) 的一个法向量为 \(\bs{n}_1=(1,0,0)\). 平面 \(PCD\) 内两条相交向量可取 \(\overrightarrow{PC}=(0,1,-2)\) 和 \(\overrightarrow{PD}=(2,0,-2)\)，故其一个法向量为 \(\bs{n}_2=(1,2,1)\). 设二面角的平面角为 \(\theta\)，则
\[
\cos\theta=\frac{|\bs{n}_1\cdot\bs{n}_2|}{|\bs{n}_1||\bs{n}_2|}=\frac1{\sqrt6}
\]
所以
\[
\sin\theta=\sqrt{1-\frac16}=\frac{\sqrt{30}}6
\]
\item 设 \(AE=h\)，则 \(E=(0,0,h)\)，且 \(0\leqslant h\leqslant2\). 取 \(\overrightarrow{BE}=\left(\frac12,-\frac12,h\right)\)，\(\overrightarrow{CD}=(2,-1,0)\).
于是
\(\overrightarrow{BE}\cdot\overrightarrow{CD}=\frac32\)，\(|BE|=\sqrt{h^2+\frac12}\)，\(|CD|=\sqrt5\).
由异面直线所成角取锐角，得
\[
\cos30^{\circ}=\frac{3/2}{\sqrt{h^2+1/2}\sqrt5}
\]
化简得 \(h^2=\frac1{10}\). 因 \(h\geqslant0\)，所以 \(AE=h=\frac{\sqrt{10}}{10}\).
\end{enumerate}
\end{solution}

\begin{problem}
已知 \(\{a_{n}\}\) 是等差数列，其前 \(n\) 项和为 \(S_{n}\)，\(\{b_{n}\}\) 是等比数列，且 \(a_{1} = b_{1} = 2\)，\(a_{4} + b_{4} = 27\)，\(S_{4} - b_{4} = 10\).

\begin{enumerate}
\item 求数列 \( \{a_{n}\} \) 与 \( \{b_{n}\} \) 的通项公式；
\item 记 \(T_{n} = a_{n}b_{1} + a_{n - 1}b_{2} + \dots +a_{1}b_{n}\)，\(n \in \N^*\)，证明 \(T_{n} + 12 = -2a_{n} + 10b_{n}\) （\(n \in \N^*\)）.
\end{enumerate}

\end{problem}

\begin{solution}
设等差数列 \(\{a_n\}\) 的公差为 \(d\)，等比数列 \(\{b_n\}\) 的公比为 \(q\)，则 \(a_4=2+3d\)，\(b_4=2q^3\)，\(S_4=8+6d\). 由题设得
\[
\begin{cases}
2+3d+2q^3=27,\\
8+6d-2q^3=10
\end{cases}
\]
两式相加得 \(10+9d=37\)，所以 \(d=3\). 代入第一式得 \(2q^3=16\)，从而 \(q=2\). 因此 \(a_n=2+3(n-1)=3n-1\)，\(b_n=2\cdot2^{n-1}=2^n\).
\begin{enumerate}
\item 由上面的计算，\(a_n=3n-1\)，\(b_n=2^n\).
\item 由 \(b_{n+1}=2b_n\)，有
\[
T_{n+1}=a_{n+1}b_1+a_nb_2+\cdots+a_1b_{n+1}
\]
而
\[
2T_n=a_nb_2+a_{n-1}b_3+\cdots+a_1b_{n+1}
\]
所以 \(T_{n+1}=2T_n+2a_{n+1}=2T_n+6n+4\). 当 \(n=1\) 时，\(T_1=a_1b_1=4\)，且
\[
T_1+12=16=-2a_1+10b_1
\]
假设对某个正整数 \(n\) 有 \(T_n+12=-2a_n+10b_n\)，则
\[
T_{n+1}+12=2T_n+6n+16=2(-2a_n+10b_n-12)+6n+16
\]
从而
\[
T_{n+1}+12=-4(3n-1)+20b_n+6n-8=-6n-4+20b_n=-2a_{n+1}+10b_{n+1}
\]
其中使用了 \(a_{n+1}=3(n+1)-1\) 和 \(b_{n+1}=2b_n\). 由数学归纳法，对任意 \(n\in\N^*\)，都有 \(T_n+12=-2a_n+10b_n\).
\end{enumerate}
\end{solution}

\begin{problem}
设椭圆 \(\frac{x^2}{a^2} + \frac{y^2}{b^2} = 1\)（\(a > b > 0\)）的左，右顶点分别为 \(A\)，\(B\)，点 \(P\) 在椭圆上且异于 \(A\)，\(B\) 两点，点 \(O\) 为坐标原点.

\begin{enumerate}
\item 若直线 \(AP\) 与 \(BP\) 的斜率之积为 \(-\frac{1}{2}\)，求椭圆的离心率；
\item 若 \( \left|AP\right|=\left|OA\right| \)，证明直线 \(OP\) 的斜率 \(k\) 满足 \( \left|k\right|>\sqrt{3} \).
\end{enumerate}

\end{problem}

\begin{solution}
设 \(P=(x_0,y_0)\). 因 \(P\) 不是左右顶点，故 \(y_0\neq0\).
\begin{enumerate}
\item 椭圆左右顶点为 \(A=(-a,0)\)，\(B=(a,0)\)，所以 \(k_{AP}=\frac{y_0}{x_0+a}\)，\(k_{BP}=\frac{y_0}{x_0-a}\).
由椭圆方程，\(y_0^2=\frac{b^2}{a^2}(a^2-x_0^2)\)，故
\[
k_{AP}k_{BP}=\frac{y_0^2}{x_0^2-a^2}=-\frac{b^2}{a^2}
\]
由题设 \(-\frac{b^2}{a^2}=-\frac12\)，所以 \(\frac{b^2}{a^2}=\frac12\). 椭圆离心率为
\[
e=\frac{\sqrt{a^2-b^2}}a=\sqrt{1-\frac{b^2}{a^2}}=\frac{\sqrt2}{2}
\]
\item 令 \(r=\frac{b^2}{a^2}\)，则 \(0<r<1\)，再令 \(u=\frac{x_0}{a}\). 由 \(|AP|=|OA|=a\) 和椭圆方程，得
\[
(u+1)^2+r(1-u^2)=1
\]
即 \((1-r)u^2+2u+r=0\). 在 \([-1,1]\) 上，该式左端的导数为 \(2(1-r)u+2\geqslant2r>0\)，故左端严格递增；又
\[
(1-r)\left(-\frac12\right)^2+2\left(-\frac12\right)+r=\frac{3(r-1)}4<0
\]
所以方程在 \((-\frac12,0)\) 内的根就是 \(u\). 由于 \(u\neq0\)，直线 \(OP\) 的斜率存在，且
\[
k^2=\frac{y_0^2}{x_0^2}=\frac{r(1-u^2)}{u^2}=-1-\frac2u
\]
从而
\[
k^2-3=-1-\frac2u-3=\frac{-2(2u+1)}u>0
\]
因为 \(-\frac12<u<0\). 故 \(|k|>\sqrt3\).
\end{enumerate}
\end{solution}

\begin{problem}
已知函数 \( f(x) = x - \ln (x + a) \) 的最小值为 0，其中 \( a > 0 \).

\begin{enumerate}
\item 求 \(a\) 的值；
\item 若对任意的 \(x \in [0, +\infty)\)，有 \(f(x) \leqslant kx^2\) 成立. 求实数 \(k\) 的最小值；
\item 证明：\(\displaystyle\sum_{i=1}^{n} \frac{2}{2i-1} - \ln(2n+1) < 2\)（\(n \in \N^*\)）.
\end{enumerate}

\end{problem}

\begin{solution}
函数定义域为 \((-a,+\infty)\)，且
\[
f'(x)=1-\frac1{x+a}=\frac{x+a-1}{x+a}
\]
因此当 \(x\in(-a,1-a)\) 时 \(f'(x)<0\)，当 \(x\in(1-a,+\infty)\) 时 \(f'(x)>0\)，函数在 \(x=1-a\) 处取得最小值. 由最小值为零，得
\[
f(1-a)=1-a-\ln1=1-a=0
\]
所以 \(a=1\).
\begin{enumerate}
\item \(a=1\)，故 \(f(x)=x-\ln(1+x)\).
\item 若 \(k\leqslant0\)，取 \(x=1\)，有 \(f(1)=1-\ln2>0\)，而 \(kx^2\leqslant0\)，不满足题设. 因此只需考察 \(k>0\). 令
\[
h(x)=\frac{x^2}{2}-x+\ln(1+x)
\]
则 \(h(0)=0\)，且
\[
h'(x)=x-1+\frac1{1+x}=\frac{x^2}{1+x}\geqslant0
\]
所以 \(h(x)\geqslant0\)，即 \(x-\ln(1+x)\leqslant\frac{x^2}{2}\)，故 \(k=\frac12\) 时不等式对任意 \(x\geqslant0\) 成立. 另一方面，当 \(x>0\) 时，任何满足条件的 \(k\) 都有
\[
k\geqslant\frac{x-\ln(1+x)}{x^2}
\]
而
\[
\lim_{x\to0^+}\frac{x-\ln(1+x)}{x^2}=\lim_{x\to0^+}\frac{1-1/(1+x)}{2x}=\frac12
\]
所以 \(k\geqslant\frac12\). 因此 \(k\) 的最小值为 \(\frac12\).
\item 当 \(n=1\) 时，左端为 \(2-\ln3<2\). 当 \(n\geqslant2\) 时，令 \(u=\frac1{i-1}>0\)，则
\[
\ln\frac{i}{i-1}-\frac2{2i-1}=\ln(1+u)-\frac{2u}{2+u}
\]
记右端为 \(r(u)\)，有 \(r(0)=0\)，且
\[
r'(u)=\frac1{1+u}-\frac4{(2+u)^2}=\frac{u^2}{(1+u)(2+u)^2}>0
\]
所以 \(\ln\frac{i}{i-1}>\frac2{2i-1}\). 于是
\[
\sum_{i=1}^{n}\frac2{2i-1}=2+\sum_{i=2}^{n}\frac2{2i-1}<2+\sum_{i=2}^{n}\ln\frac{i}{i-1}=2+\ln n<2+\ln(2n+1)
\]
移项即得 \(\displaystyle\sum_{i=1}^{n}\frac2{2i-1}-\ln(2n+1)<2\).
\end{enumerate}
\end{solution}
